再次强调，这通常被不正式地写成：
\begin{equation}
  \delta CS(A)=\frac{1}{8\pi^{2}}\left[2\operatorname{tr}\delta A\,F-d\left(\operatorname{tr}(A\,\delta A)\right)\right]
  \tag{3.14}
  \label{eq:3.14}
\end{equation}
% [Tutor Note: (3.14) 的归一化极易与前面阿贝尔情形的约定混淆。(3.14) 与 (1.6) 同为「无 $1/2$、带 $1/8\pi^{2}$」的非阿贝尔写法，而第 2 章 $U(1)$ 情形用的是 $1/4\pi$ 系数且 $A$ 取实值（对比 (3.15) 的 $k/4\pi$ 与 \eqref{eq:2.79}）。把 (3.14) 用到具体计算之前，先确认三件事：$\operatorname{tr}$ 是否取基础表示、$d$ 是否作用在 $A\,\delta A$ 上、以及全导数项在 $\partial M_3=\varnothing$ 时才可丢弃。符号上 $2\operatorname{tr}\delta A\,F$ 的因子 2 来自 $\delta(A\,dA)$ 的两项合并，丢掉它会连带丢掉 (3.16) 中二次项的因子 2。]

{\small\itshape ［原文标记：需要再次评注非阿贝尔情形下的旋量理论。］}
% [原文待办标记，原文 p.88；原文以两行星号线夹住，为作者提醒自己补写的内容]

\subsubsection{\texorpdfstring{$M_3$}{M3} 具有非空边界时的 Chern-Simons 作用量} % 原文: 3.1.1 The Chern-Simons Action On M3 With A Nonempty Boundary

{\small\itshape ［原文标记：要说明此时作用量必须被看作某个线丛的一个截面。］}
% [原文待办标记，原文 p.88]

\subsection{Chern-Simons 路径积分的半经典近似} % 原文: 3.2 The Semiclassical Approximation To The Chern-Simons Path Integral

于是我们现在考虑把 (3.15) 推广而成的形式路径积分：
% 转写说明：原文此处写 generalizing (3.15)，而紧随其后的公式正是编号 (3.15)，显系自指，疑为对前面某式的笔误。照原文保留编号不动，不改为 \eqref{eq:3.15}。
\begin{equation}
  \int_{\mathcal{A}/\mathcal{G}}[dA]\,\exp\left[\mathrm{i}\frac{k}{4\pi}\int_{M_3}\operatorname{tr}\left(A\,dA+\frac{2}{3}A^{3}\right)\right]\prod_{\alpha}W(R_{\alpha},\gamma_{\alpha})
  \tag{3.15}
  \label{eq:3.15}
\end{equation}

眼下我们取 $\gamma_{\alpha}$ 为闭曲线，并取 $M_3$ 没有边界。

形式上我们预期这是一个拓扑场论，因此它应当只依赖于 $M_3$ 的拓扑以及链环 $\bigsqcup_{\alpha}\gamma_{\alpha}$ 的拓扑。

当然，由我们对 $U(1)$ 情形的经验，我们预期拓扑不变性会出现一些反常。

在半经典近似 $k\to\infty$ 下，我们应当在运动方程的解附近展开。这些解就是 $M_3$ 上的平坦联络。与 $\Sigma_2\times S^{1}$ 这类简单情形形成鲜明对照的是，实际上此时往往只有孤立的平坦联络。

{\small\itshape ［原文标记：一行星号分隔线。］}
% [原文待办标记，原文 p.88；作者以星号行分隔草稿段落，其下即为 "Examples" 小标题]

\textbf{例子}

\begin{enumerate}
  \item \emph{有理同调球}（rational homology spheres）：取 $M_3=S^{3}/\Gamma$，其中 $\Gamma$ 是 $SU(2)$ 的离散子群。则 $\pi_1(M_3)\cong\Gamma$。若 $G=SU(2)$，则（在共轭意义下）这样的子群唯一，因此平坦联络是孤立的。

  \item \emph{映射环面}（mapping torus）。设 $S$ 是一个曲面，$f:S\to S$ 是一个微分同胚。则可构造映射环面 $M_f=(S\times[0,1])/\sim$，其中 $(x,0)\sim(f(x),1)$。它可以看作以 Riemann 曲面为纤维、以圆为底的纤维化。{\small\itshape ［原文标记：关于双曲结构的 Agol 定理。此处需要对这一情形的基本群作些说明。］}
  % [原文待办标记，原文 p.88]

  \item \emph{Heegaard 分裂}（Heegaard splittings）。
\end{enumerate}

\subsubsection{半经典近似：再论标架反常} % 原文: 3.2.1 Semiclassical Approximation: The Framing Anomaly Revisited

我们现在考虑 $k\to\infty$ 的渐近行为。由于 $\hbar\sim1/k$，这就是半经典近似。我们遵从 Witten 的论文 [67] 中的讨论，只作若干修改。

我们应当在运动方程的解附近展开。这些解是平坦联络。为简单起见，我们假定这些平坦联络（在规范等价意义下）是孤立的。

于是让我们考虑在一个孤立的平坦联络 $A_{*}$ 附近展开。我们写 $A=A_{*}+\alpha$，并利用式 \eqref{eq:3.6} 得到作用量为
% 跨片段引用说明：\eqref{eq:3.6} 指向 c10 片段中的 (3.6)（Chern-Simons 形式的一阶变分），其 \label 由该片段提供，本片段不重复引入。
\begin{equation}
  \frac{k}{4\pi}\int_{M_3}\operatorname{tr}\left(A\,dA+\frac{2}{3}A^{3}\right)
  =\frac{k}{4\pi}\int_{M_3}\operatorname{tr}\left(A_{*}\,dA_{*}+\frac{2}{3}A_{*}^{3}\right)
  +\frac{k}{4\pi}\int\operatorname{tr}\left(\alpha\,D_{A_{*}}\alpha+\frac{2}{3}\alpha^{3}\right)
  \tag{3.16}
  \label{eq:3.16}
\end{equation}
% [Tutor Note: (3.16) 中间省略了一步。用 $\mathrm{cs}(A_{*}+\alpha)=\mathrm{cs}(A_{*})+2\operatorname{tr}(\alpha F_{A_{*}})+\operatorname{tr}(\alpha D_{A_{*}}\alpha)+\frac{2}{3}\operatorname{tr}\alpha^{3}$ 逐项比对即可：一次项是 $\frac{k}{2\pi}\int\operatorname{tr}(\alpha F_{A_{*}})$，正是「$A_{*}$ 平坦」把它整个消掉，才有「经典作用量 + 纯二次 + 纯三次」的干净结构（$A_{*}$ 不平坦时此式还缺一个线性项）；而 $\alpha\,D_{A_{*}}\alpha$ 是把 $\alpha\,d\alpha+\alpha[A_{*},\alpha]$ 合写成协变外微分的产物。另有两处照原文保留未改：右端第三个积分漏印了下标 $M_3$（应按 $\int_{M_3}$ 读）；$A_{*}^{3}$ 前的系数 $\frac{2}{3}$ 与 $\alpha^{3}$ 的同源，不要因为两者相同就以为写重了。]

右端第一项是经典作用量 $S(A_{*})$。于是我们得到一个非退化的二次项和一个相互作用项。为作单圈近似，我们完全如上面的方程 {\small\itshape ［原文标记：＊＊＊，指代待补的公式编号］} 那样引入 BRST 鬼场，并利用规范固定条件
\begin{equation}
  D_{A}*\alpha=D^{i}_{A}\alpha_{i}\operatorname{vol}(g)=0
  \tag{3.17}
  \label{eq:3.17}
\end{equation}
引入一个度规，从而得到二次作用量：
\begin{equation}
  \exp\left[\mathrm{i}\frac{k}{4\pi}\int_{M_3}\operatorname{tr}\left(\alpha\,D_{A_{*}}\alpha+H\,D_{A_{*}}\alpha+b\,D_{A_{*}}*D_{A_{*}}c\right)\right]
  \tag{3.18}
  \label{eq:3.18}
\end{equation}
这里 $\alpha\in\Omega^{1}(M_3;\mathrm{ad}P)$，而 $H,b,c\in\Omega^{0}(M_3;\mathrm{ad}P)$，其中 $b,c$ 是反交换场，而 $H$ 是交换场。

现在，由于 $A_{*}$ 是平坦联络，$D_{A_{*}}:\Omega^{k}(M_3;\mathrm{ad}P)\to\Omega^{k+1}(M_3;\mathrm{ad}P)$ 的平方为零：
\begin{equation}
  (D_{A_{*}})^{2}=0
  \tag{3.19}
  \label{eq:3.19}
\end{equation}
因此 $D_{A_{*}}$ 充当微分。于是情形与上面第 2.2.17 节所讨论的阿贝尔情形极为类似。附录 F 的 Hodge 理论讨论以同样的方式适用，于是我们得到单圈行列式
\begin{equation}
  \frac{\det'(*D_{A_{*}}*D_{A_{*}})}{\det^{1/2}\mathcal{D}}
  \tag{3.20}
  \label{eq:3.20}
\end{equation}
其中
\begin{equation}
  \mathcal{D}:\Omega^{\mathrm{odd}}(M_3;\mathrm{ad}P)\to\Omega^{\mathrm{odd}}(M_3;\mathrm{ad}P)
  \tag{3.21}
  \label{eq:3.21}
\end{equation}
由 $\mathcal{D}=*D_{A_{*}}+D_{A_{*}}*$ 定义。附录 F 的 Hodge 理论论证表明 \eqref{eq:3.20} 的绝对值是拓扑不变量，但其相位原来是连续依赖于度规的。

{\small\itshape ［原文标记：某处需要引用 Albert Schwarz 的工作。］}
% [原文待办标记，原文 p.89，位于页边]

为了研究度规依赖性，我们把 $\mathcal{D}$ 解释为一个扭化 Dirac 算符。回想反对称张量可以看作双旋量：
\begin{equation}
  S\otimes S\cong\oplus_{k}\Lambda^{k}T^{*}M_3
  \tag{3.22}
  \label{eq:3.22}
\end{equation}

在这一同构下，耦合到旋量的 Dirac 算符变成作用在微分形式上的第一阶微分算符，事实上变成 $d*+*d$。若耦合到 $\mathrm{ad}P$，情况也一样。

现在，三维 Dirac 算符的相位有一个一般公式 [4, 72]
\begin{equation}
  \operatorname{Det}\gamma\cdot\mathcal{D}=|\operatorname{Det}\gamma\cdot\mathcal{D}|\,e^{\frac{\mathrm{i}\pi}{2}\eta(\gamma\cdot\mathcal{D})}
  \tag{3.23}
  \label{eq:3.23}
\end{equation}

{\small\itshape ［原文标记：需要定义 $\eta$ 不变量；解释 [4] 中的论证；也许要在更前面解释 Fermi 行列式的相位与 $\eta$？］}
% [原文待办标记，原文 p.90；原文以两行星号线夹住]

$\eta$ 不变量 $\eta(\mathcal{D})$ 同时依赖于平坦联络 $A_{*}$ 和度规。对 $A_{*}$ 的依赖可以通过考虑从 $A_{*}$ 到平凡联络 $A=0$ 的一族单参数联络 $A(s)$ 来算出。（由于我们讨论的是平坦联络，丛 $P$ 可平凡化。）{\small\itshape ［原文标记：且慢！只有当 $G$ 单连通时这才是对的。这里有一个漏洞！！］}
% [原文待办标记，原文 p.90，位于页边]

然后我们把 APS 指标定理应用于四维流形 $M_3\times[0,1]$，其中 $s$ 是沿区间的坐标。把 $\eta(\mathcal{D})$ 记作 $\eta(A_{*},g_{\mu\nu})$ 以强调对这些变量的依赖，由 APS 定理得到：{\small\itshape ［原文标记：这里需要补上一些细节！！！］}
% [原文待办标记，原文 p.90，位于页边]
\begin{equation}
  \frac{\pi}{2}\eta(A_{*},g_{\mu\nu})-\frac{\pi}{2}\eta(0,g_{\mu\nu})=\frac{1}{2}c_{2}(G)\,2\pi\int_{M_3}CS(A_{*})
  \tag{3.24}
  \label{eq:3.24}
\end{equation}

右端再次是一个拓扑不变量。然而，此时我们必须面对 $\eta(0,g_{\mu\nu})$ 的度规依赖性。于是我们考虑一族单参数度规 $g_{\mu\nu}(s)dx^{\mu}\otimes dx^{\nu}$，并再次把 APS 指标定理应用于四维流形 $M_3\times[0,1]$，其中 $s$ 是沿区间的坐标。现在，当我们把规范场取为 $A=0$ 后，规范指标已经解耦，于是
\begin{equation}
  \eta(0,g_{\mu\nu})=(\dim G)\,\eta(g_{\mu\nu})
  \tag{3.25}
  \label{eq:3.25}
\end{equation}
其中 $\eta(g_{\mu\nu})$ 是算符 $*d+d*$ 限制在奇形式上的 $\eta$ 不变量。现在 APS 指标定理给出：
\begin{equation}
  \operatorname{Index}\gamma\cdot\mathcal{D}=\int_{M_3\times[0,1]}\hat{A}-\frac{1}{2}\left[\eta(g^{(1)}_{\mu\nu})-\eta(g^{(1)}_{\mu\nu})\right]
  \tag{3.26}
  \label{eq:3.26}
\end{equation}
% [Tutor Note: (3.26) 方括号内两个 $\eta$ 的上标原文都印作 $(1)$，显系笔误——应是一端取 $s=0$ 的度规、另一端取 $s=1$ 的度规，即 $\eta(g^{(1)}_{\mu\nu})-\eta(g^{(0)}_{\mu\nu})$；否则括号恒为零而 $\operatorname{Index}$ 与 $\hat{A}$ 积分一般不为零，等式不成立。同理，(3.27) 的两个右端逐字相同，也是重复排印。两处均照原文转写未改；实际要用的是四维流形上 $\hat{A}$ 展开只有六次项 $\hat{A}=-\frac{1}{24}p_{1}(R)$，再配 (3.28) 的 Chern-Weil 代表元。]
但是
\begin{equation}
  \int_{M_3\times[0,1]}\hat{A}=-\frac{1}{24}\int_{M_3\times[0,1]}p_{1}(R)=-\frac{1}{24}\int_{M_3\times[0,1]}p_{1}(R)
  \tag{3.27}
  \label{eq:3.27}
\end{equation}
其中
\begin{equation}
  p_{1}(R)=-\frac{1}{8\pi^{2}}\operatorname{Tr}_{\mathrm{vector}}R^{2}
  \tag{3.28}
  \label{eq:3.28}
\end{equation}
是 Pontryagin 类的 Chern-Weil 代表元。于是
\begin{equation}
  \int_{[0,1]}\operatorname{Tr}_{\mathrm{vector}}R^{2}=CS(\omega^{1},\omega^{2})
  \tag{3.29}
  \label{eq:3.29}
\end{equation}

其中 $\omega$ 是 Levi-Civita 联络的旋量联络（spin connection），故 $R=d\omega+\omega^{2}$。非正式地，
\begin{equation}
  CS(\omega)=\operatorname{Tr}_{\mathrm{vector}}\left(\omega\,d\omega+\frac{2}{3}\omega^{3}\right)
  \tag{3.30}
  \label{eq:3.30}
\end{equation}
由此可知，如果我们平凡化切丛、从而可以使 $CS(\omega)$ 具有整体意义，那么
\begin{equation}
  \mathcal{F}:=\frac{1}{2}\eta(g_{\mu\nu})+\frac{1}{24}\int_{M_3}\frac{1}{8\pi^{2}}\operatorname{Tr}_{\mathrm{vector}}\left(\omega\,d\omega+\frac{2}{3}\omega^{3}\right)
  \tag{3.31}
  \label{eq:3.31}
\end{equation}
{\small\itshape ［原文标记：需要说明究竟在多大程度上确实可以平凡化切丛。］}
% [原文待办标记，原文 p.91，位于页边]
与度规无关。但 $\mathcal{F}$ 依赖于切丛的标架（framing）。给定一个标架以及一个映射 $\rho:M_3\to SO(3)$，我们可以造出另一个标架。旋量联络的 Chern-Simons 不变量会改变 $\rho$ 的绕数 $w(\rho)$ 的若干倍，于是
\begin{equation}
  \mathcal{F}\to\mathcal{F}+\frac{w(\rho)}{12}
  \tag{3.32}
  \label{eq:3.32}
\end{equation}

\textbf{评论}

\begin{enumerate}
  \item Witten 论证说，我们应当加上一个正比于引力 Chern-Simons 作用量的局域抵消项，以消去度规依赖。于是单圈近似成为：
\begin{equation}
  Z\sim_{k\to\infty}e^{\mathrm{i}\pi\mathcal{F}}\sum_{A_{*}}T(A_{*})\,e^{2\pi\mathrm{i}\left(k+\frac{1}{2}C_{2}(G)\right)\int_{M_3}CS(A_{*})}
  \tag{3.33}
  \label{eq:3.33}
\end{equation}
  其中 $T(A_{*})$ 是 $D_{A_{*}}$ 的解析挠率（analytic torsion）。

  \item 相位移动为
\begin{equation}
  e^{\mathrm{i}\pi\mathcal{F}}\to e^{2\pi\mathrm{i}\frac{d}{24}w(\rho)}\,e^{\mathrm{i}\pi\mathcal{F}}
  \tag{3.34}
  \label{eq:3.34}
\end{equation}
  现在，由下文所述的与 RCFT 的关系，我们其实预期这一移动应当是
\begin{equation}
  e^{\mathrm{i}\pi\mathcal{F}}\to e^{2\pi\mathrm{i}\frac{c}{24}w(\rho)}\,e^{\mathrm{i}\pi\mathcal{F}}
  \tag{3.35}
  \label{eq:3.35}
\end{equation}
  其中 $c$ 是 WZW 理论的中心荷：
\begin{equation}
  c=\frac{k}{k+h^{\vee}}\dim G
  \tag{3.36}
  \label{eq:3.36}
\end{equation}
  由于我们看的是首阶 $k\to\infty$ 渐近，这并不立即构成矛盾，但它确实提出了 $1/k$ 修正项的度规依赖性问题。
% [Tutor Note: (3.34) 与 (3.35) 的「不矛盾」正是这段的要点，值得自己验一遍：把 (3.36) 在 $k\to\infty$ 展开，$c=\dim G\,(1-h^{\vee}/k+\cdots)\to\dim G=:d$，故 $c/24$ 与 $d/24$ 之差是 $\mathcal{O}(1/k)$，正落在 (3.33) 首阶单圈精度之外。作者真正的担心是：若把 (3.33) 推进到 $1/k$ 阶，$\mathcal{F}$ 的标架依赖就会以 $c$ 而非 $d$ 出现，那时 (3.31) 的抵消项选择需重新审视。]

  \item 因此，路径积分所依赖的拓扑数据的选取包括：定向的选取与标架的选取。当我们纳入 Wilson 线时，我们还需要把圈的标架一并纳入，正如我们已在阿贝尔情形所见。
\end{enumerate}

{\small\itshape ［原文标记：一行星号分隔线。］}
% [原文待办标记，原文 p.91；其下为作者草拟的补充文献清单 "Remarks from"]

\textbf{补注（Remarks from）}

\begin{enumerate}
  \item Atiyah 关于 2-framing 的论文。

  \item Segal 关于 rigging 的工作。

  \item $p_{1}$ 结构：Blanchet、Habegger、Masbaum、Vogel。
\end{enumerate}

{\small\itshape ［原文标记：一行星号分隔线。］}
% [原文待办标记，原文 p.92；作者以此结束草稿性质的补注清单]

\subsubsection{高阶圈} % 原文: 3.2.2 Higher Loops

评注：Axelrod-Singer。

\subsection{若干重要情形下的量子化} % 原文: 3.3 Quantization In Some Important Cases

\subsubsection{\texorpdfstring{$\Sigma_2\times\mathbb{R}$}{Σ2×R} 上的量子化：\texorpdfstring{$\Sigma$}{Σ} 为紧致曲面} % 原文: 3.3.1 Quantization On Σ2 × R With Σ A Compact Surface

我们现在对 $M_3=\Sigma_2\times\mathbb{R}$ 的情形量子化该理论，其中 $\Sigma_2$ 是一个定向曲面。因此我们假定有一个主 $G$-丛 $P\to\Sigma_2$，并先量子化全部规范场构成的空间 $\mathcal{A}=\mathrm{Conn}(P\to\Sigma_2)$，然后强加 Gauss 定律。

回想 $\mathcal{A}$ 上有辛形式
\begin{equation}
  \omega=\frac{k}{4\pi}\int_{\Sigma}\operatorname{tr}\left(\delta A_{s}\wedge\delta A_{s}\right)
  \tag{3.37}
  \label{eq:3.37}
\end{equation}
它给出 Poisson 括号：
\begin{equation}
  \{A^{a}_{i}(x),A^{b}_{j}(y)\}=\frac{2\pi}{k}\epsilon_{ij}\delta^{ab}\delta^{(2)}(x,y)
  \tag{3.38}
  \label{eq:3.38}
\end{equation}
其中 $A=A^{a}_{i}t_{a}dx^{i}$，而 $t^{a}$ 是关于迹 $\operatorname{tr}$ 的一组正交归一基。规范变换 $\epsilon:\Sigma_2\to\mathfrak{g}$ 的动量映射为
\begin{equation}
  \mu(\epsilon)=-\frac{k}{2\pi}\int_{\Sigma}\operatorname{tr}\left(\epsilon F\right)
  \tag{3.39}
  \label{eq:3.39}
\end{equation}

现在我们需要量子化。对一切曲面 $\Sigma_2$ 都可用的一种办法是选取一个复结构，使之成为 Riemann 曲面。局部地，我们可以定义全纯坐标 $z=x^{1}+\mathrm{i}x^{2}$，其中 $dx^{1}dx^{2}$ 为正向的；同时还选取丛的一个局部平凡化，并定义 $A=A_{z}dz+A_{\bar z}d\bar z$，于是得到
\begin{equation}
  [A^{a}_{z}(x),A^{b}_{\bar z}(y)]=\frac{\pi}{k}\delta^{ab}\delta^{(2)}(x-y)
  \tag{3.40}
  \label{eq:3.40}
\end{equation}
这在空间 $\mathcal{A}$ 上诱导出一个复结构：
\begin{equation}
  T^{1,0}\mathcal{A}\cong\Omega^{1,0}(\Sigma_2;\mathrm{ad}P)
  \tag{3.41}
  \label{eq:3.41}
\end{equation}
此外，我们可以用相容的度规把 $\mathcal{A}$ 变成一个 Kähler 流形：
\begin{equation}
  g(\alpha,\alpha)=\frac{k}{4\pi}\int_{\Sigma_2}\operatorname{tr}\left(\alpha^{1,0}\alpha^{0,1}\right)
  \tag{3.42}
  \label{eq:3.42}
\end{equation}
其中 $\alpha\in\Omega^{1}(\Sigma_2;\mathrm{ad}P)$ 是整体定义的、取值于伴随丛的 1-形式。预量子线丛 $\mathcal{L}_{k}\to\mathcal{A}$ 具有 Kähler 势：
\begin{equation}
  K=\frac{k}{4\pi}\int_{\Sigma_2}\operatorname{tr}\left(A^{1,0}A^{0,1}\right).
  \tag{3.43}
  \label{eq:3.43}
\end{equation}
% [Tutor Note: (3.38) 到 (3.40) 的系数变化（$2\pi/k\to\pi/k$）务必自己算一遍，否则后面的 $k/\pi$ 与 $\xi$ 会一路错。代入 $A_{z}=\frac{1}{2}(A_{1}-\mathrm{i}A_{2})$、$A_{\bar z}=\frac{1}{2}(A_{1}+\mathrm{i}A_{2})$，并与 (3.38) 的 $\epsilon_{12}=+1$ 约定一致，直接组合得到的是 $\frac{\mathrm{i}\pi}{k}\delta^{ab}\delta^{(2)}(x-y)$，即右端带一个 $\mathrm{i}$；(3.40) 写成实系数 $\pi/k$，说明作者已把几何量子化中 $\{\cdot,\cdot\}\to\frac{1}{\mathrm{i}}[\cdot,\cdot]$ 的 $\mathrm{i}$ 吸收进约定（与 (3.44)–(3.47) 里 $\mu(\epsilon)\Psi=0$、$F\Psi=0$ 的算符化方式配套）。核对 (3.47) 的 $k/\pi$ 时必须沿用同一约定。]

现在规范变换辛地作用，我们可以尝试作 Kähler 量子化。若在经典层面强加约束，我们就被引向平坦规范场模空间的 Kähler 量子化。这会很快把人带入代数几何的世界。参见 3.4 的简要评注。我们转而将在量子层面强加约束。

若 $\Psi\in\Gamma[\mathcal{L}_{k}\to\mathcal{A}]$，我们可以形式上把它看作全纯函数 $\Psi[A_{z}]$。我们将需要强加 Gauss 定律
\begin{equation}
  \mu(\epsilon)\Psi=0
  \tag{3.44}
  \label{eq:3.44}
\end{equation}
并利用量子化规则 (3.40) 得到一个关于 $A_{z}$ 的第一阶（泛函）微分方程。我们把
\begin{equation}
  \frac{\delta}{\delta A_{z}(x)}\Psi\in T^{1,0}_{x}\otimes\mathfrak{g}\otimes\mathcal{L}_{k,A}
  \tag{3.45}
  \label{eq:3.45}
\end{equation}
那么算符方程
\begin{equation}
  F\Psi=0
  \tag{3.46}
  \label{eq:3.46}
\end{equation}
就变成第一阶微分方程
\begin{equation}
  D_{A,z}\frac{\delta}{\delta A_{z}(x)}\Psi=-\frac{k}{\pi}\bar\partial_{\bar z}A_{z}(x)\Psi
  \tag{3.47}
  \label{eq:3.47}
\end{equation}
% [Tutor Note: (3.46) 到 (3.47) 之间跳了一步，补出来是：$\Sigma_2\times\mathbb{R}$ 上 $F_{z\bar z}=0$ 的分量式即 $\partial_{\bar z}A_{z}-\partial_{z}A_{\bar z}+[A_{\bar z},A_{z}]=0$，而按 (3.40) 有 $A_{\bar z}\Psi=-\frac{\pi}{k}\frac{\delta}{\delta A_{z}}\Psi$（符号由 $\{\cdot,\cdot\}\to\frac{1}{\mathrm{i}}[\cdot,\cdot]$ 与 $\bar\partial$ 的取向共同决定），代回即得 (3.47)。核对时重点看三处：$k/\pi$ 的倒数关系、$\bar\partial_{\bar z}$ 只作用在 $A_{z}$ 上、以及 $D_{A,z}$ 是协变 $\partial_{z}$（含 $[A_{z},\cdot]$）。]

正如我们将看到的，这个微分方程将要求 $\Psi$ 位于一个有限维向量空间中。我们还将把它解释为关于流守恒中反常的一个方程。{\small\itshape ［原文标记：需要回头解释为什么这是 WZW 模型中整体对称性反常的一个陈述。］}
% [原文待办标记，原文 p.93，位于页边]

在「$\mathcal{L}_{k}$ 的 $L^{2}$ 截面」构成的 Hilbert 空间上，我们将有规范变换群 $\mathrm{Aut}(P)$ 的酉算符表示。当 $P$ 可平凡化时，这只是由映射 $g:\Sigma_2\to G$ 构成的群（须满足适当的正则性条件）。我们想要强加
\begin{equation}
  U(g)\Psi=\Psi
  \tag{3.48}
  \label{eq:3.48}
\end{equation}
其中
\begin{equation}
  (U(g)\cdot\Psi)(A_{z})=e^{f(A_{z},g)}\,\Psi((A^{g})_{z})
  \tag{3.49}
  \label{eq:3.49}
\end{equation}
注意，由于 $d+A^{g}=g^{-1}(d+A)g$，而这是联络空间上的右作用，我们在截面 Hilbert 空间上得到的是左作用。

先验地，可能存在射影表示：
\begin{equation}
  U(g_{1})U(g_{2})=c(g_{1},g_{2})\,U(g_{1}g_{2})
  \tag{3.50}
  \label{eq:3.50}
\end{equation}
其中 $c(g_{1},g_{2})$ 是某个非平凡上循环；然而我们想要强加 \eqref{eq:3.48}，因此该上循环应当是平凡的。

现在，微分方程 \eqref{eq:3.47} 清楚地表明，在 \eqref{eq:3.49} 中不能取 $f(A_{z};g)=0$，因为我们可以用它来计算无穷小规范变分：{\small\itshape ［原文标记：这里有一个由上述约定固定的归一化因子 $\xi$。把它算出来。］}
% [原文待办标记，原文 p.93，位于页边]
\begin{align}
  \Psi[A_{z}+D_{z}\epsilon]-\Psi[A_{z}]
  &=\xi\int_{\Sigma_2}\operatorname{tr}\left(D_{A,z}\frac{\delta}{\delta A_{z}(x)}\Psi\right)\frac{\mathrm{i}}{2}dz\,d\bar z\notag\\
  &=-\xi\frac{k}{\pi}\int_{\Sigma_2}\operatorname{tr}\left(\epsilon\,\bar\partial_{\bar z}A_{z}\right)\frac{\mathrm{i}}{2}dz\,d\bar z\,\Psi\notag\\
  &=\xi\frac{k}{\pi}\int_{\Sigma_2}\operatorname{tr}\left(\bar\partial_{\bar z}\epsilon\,A_{z}\right)\frac{\mathrm{i}}{2}dz\,d\bar z\,\Psi
  \tag{3.51}
  \label{eq:3.51}
\end{align}

假定 $\operatorname{tr}(A)=0$ 且 $g=e^{\epsilon}$，我们有：
\begin{equation}
  \Psi[A^{g}_{z}]=\exp\left[\frac{\mathrm{i}\xi}{2}\frac{k}{\pi}\int_{\Sigma_2}\operatorname{tr}\left(A^{1,0}g^{-1}\bar\partial g\right)\right]\Psi[A_{z}]
  \tag{3.52}
  \label{eq:3.52}
\end{equation}
尽管在这个阶上我们同样可以写成
\begin{equation}
  \Psi[A^{g}_{z}]=\exp\left[\frac{\mathrm{i}\xi}{2}\frac{k}{\pi}\int_{\Sigma_2}\operatorname{tr}\left(A^{1,0}\bar\partial g\,g^{-1}\right)\right]\Psi[A_{z}]
  \tag{3.53}
  \label{eq:3.53}
\end{equation}
% [Tutor Note: (3.52) 与 (3.53) 之所以在此阶等价，是因为 $g=e^{\epsilon}$ 时 $g^{-1}\bar\partial g=\bar\partial\epsilon+\mathcal{O}(\epsilon^{2})=\bar\partial g\,g^{-1}+\mathcal{O}(\epsilon^{2})$，两者只差 $\mathcal{O}(\epsilon^{2})$，而两式本身都只在此阶成立；到下一阶它们给出的 $S(g)$ 并不相同，这正是本节末尾那个 PUZZLE 的来源。另注意 (3.51)→(3.52) 用了 $\operatorname{tr}(A)=0$，把 $\int\operatorname{tr}(\bar\partial\epsilon\,A)$ 分部积分时弃掉的边界项与 $\operatorname{tr}(\epsilon\bar\partial A)$ 一并消去。]

把群律 \eqref{eq:3.50}（带平凡上循环）与群作用 \eqref{eq:3.49} 结合，给出一个相容性条件：
\begin{equation}
  f(A_{z},g_{1})+f(A^{g_{1}}_{z},g_{2})=f(A_{z},g_{1}g_{2})\bmod 2\pi\mathrm{i}\mathbb{Z}
  \tag{3.54}
  \label{eq:3.54}
\end{equation}
而且我们知道，对 $g=e^{\epsilon}$ 有
\begin{equation}
  f(A_{z},g)=\frac{\mathrm{i}\xi}{2}\frac{k}{\pi}\int_{\Sigma_2}\operatorname{tr}\left(A^{1,0}g^{-1}\bar\partial g\right)+\mathcal{O}(\epsilon^{2})
  \tag{3.55}
  \label{eq:3.55}
\end{equation}
或
\begin{equation}
  f(A_{z},g)=\frac{\mathrm{i}\xi}{2}\frac{k}{\pi}\int_{\Sigma_2}\operatorname{tr}\left(A^{1,0}\bar\partial g\,g^{-1}\right)+\mathcal{O}(\epsilon^{2})
  \tag{3.56}
  \label{eq:3.56}
\end{equation}

这些方程存在唯一的解，它关于 $A_{z}$ 是仿射线性的；也就是说，即使 $g$ 并不接近 $g=1$，我们也说：
\begin{equation}
  f(A_{z},g)=\frac{\mathrm{i}\xi}{2}\frac{k}{\pi}\int_{\Sigma_2}\operatorname{tr}\left(A^{1,0}g^{-1}\bar\partial g\right)+S(g)
  \tag{3.57}
  \label{eq:3.57}
\end{equation}

可以导出关于 $S(g)$ 的一个泛函方程，进而（至少对位于 $\mathrm{Map}(\Sigma_2,G)$ 的单位连通分支中的 $g$）把它积分出来。其结果就是所谓的 WZW 作用量，其形式为：
\begin{equation}
  f^{+}(A,g)=-\mathrm{i}\xi k\,S^{+}(g)+\mathrm{i}\xi\frac{k}{2\pi}\int_{\Sigma_2}\operatorname{tr}\left(A^{1,0}g^{-1}\bar\partial g\right)
  \tag{3.58}
  \label{eq:3.58}
\end{equation}
\begin{equation}
  f^{-}(A,g)=\mathrm{i}\xi k\,S^{-}(g)+\mathrm{i}\xi\frac{k}{2\pi}\int_{\Sigma_2}\operatorname{tr}\left(A^{1,0}\bar\partial g\,g^{-1}\right)
  \tag{3.59}
  \label{eq:3.59}
\end{equation}
其中 $S^{\pm}$ 是 WZW 作用量：
\begin{equation}
  S^{\pm}=\frac{1}{4\pi}\int_{\Sigma_2}\operatorname{tr}\left(g^{-1}\partial g\,g^{-1}\bar\partial g\right)\pm\frac{1}{12\pi}\int_{\mathcal{B}(g)}\operatorname{tr}\left(g^{-1}dg\right)^{3}
  \tag{3.60}
  \label{eq:3.60}
\end{equation}
% [Tutor Note: 「可以导出关于 $S(g)$ 的一个泛函方程」被一句话带过，被省略的是把 (3.54) 代入 (3.57) 后得到的 $1$-上循环条件 $S(g_{1}g_{2})=S(g_{1})+S(g_{2})+\frac{\mathrm{i}\xi k}{2\pi}\int_{\Sigma_2}\operatorname{tr}(A^{1,0}\cdots)$ 型关系，其解的存在性与可积性就是 Polyakov-Wiegmann 公式（乘性泛函 $S(g_{1}g_{2})=S(g_{1})+S(g_{2})-\frac{1}{4\pi}\int\operatorname{tr}(g_{1}^{-1}\partial g_{1}\,\partial g_{2}\,g_{2}^{-1})$ 之类的变体）。要真正「积分出来」，需要这条公式以及 $\mathrm{Map}(\Sigma_2,G)$ 单位连通分支上路径的选取；文献可查 Witten [67] 或 Gawędzki 的综述。]

\eqref{eq:3.60} 中第二项就是大名鼎鼎的 WZ 项。这里 $\mathcal{B}(g)\subset G$ 是 $G$ 中的一个 3-链，使得
\begin{equation}
  \partial\mathcal{B}(g)=g(\Sigma_2)
  \tag{3.61}
  \label{eq:3.61}
\end{equation}
尽管表面上如此，这一项给出的是映射 $g:\Sigma_2\to G$ 的一个局域量子场论。当然，$\mathcal{B}(g)$ 的选取是有歧义的，但对所有紧致、单、连通且单连通的群 $G$（见 [40, 11]）：
\begin{equation}
  x_{3}=\left[\frac{1}{48\pi^{2}h^{\vee}}\operatorname{Tr}_{\mathrm{adj}}\left(g^{-1}dg\right)^{3}\right]
  \tag{3.62}
  \label{eq:3.62}
\end{equation}

上式中的 $x_{3}$ 生成 $H^{3}_{DR}(G)$ 的整系数上同调格，其中 $h^{\vee}$ 是对偶 Coxeter 数。特别地，对 $SU(2)$ 我们可以取
\begin{equation}
  x_{3}=\left[\frac{1}{24\pi^{2}}\operatorname{Tr}_{2}\left(g^{-1}dg\right)^{3}\right]
  \tag{3.63}
  \label{eq:3.63}
\end{equation}
更一般地，对 $SU(N)$ 我们可以取
\begin{equation}
  x_{3}=\left[\frac{1}{24\pi^{2}}\operatorname{Tr}_{N}\left(g^{-1}dg\right)^{3}\right]
  \tag{3.64}
  \label{eq:3.64}
\end{equation}
以生成 $H^{3}_{DR}(G)$ 中的整系数上同调格。因此，按标准的论证，在定义指数化作用量时 $\mathcal{B}(g)$ 的选取无关紧要，只要 $k$ 是整数。
% [Tutor Note: (3.62)–(3.64) 的相容性完全依赖于迹的归一化，而原文没有交代 $\operatorname{Tr}_{2}$、$\operatorname{Tr}_{N}$ 的归一化（是基础表示的迹，还是它的某个整数倍）。核对办法是令 $SU(2)$（$h^{\vee}=2$）时两式给出同一个上同调类，由此读出 $\operatorname{Tr}_{\mathrm{adj}}$ 与 $\operatorname{Tr}_{2}$ 应有的比值，再看它是否与你采用的 $\operatorname{tr}(t^{a}t^{b})=\frac{1}{2}\delta^{ab}$ 约定相符；$x_{3}$ 的归一化差一个整数倍并不破坏「$k\in\mathbb{Z}$ 时 $\exp(\mathrm{i}kS_{\mathrm{WZ}})$ 单值」这一整性结论。]

{\small\itshape ［原文标记：我们的约定最好给出 $\xi=\pm1$。］}
% [原文待办标记，原文 p.95，位于页边]

最后，由上述内容可知，当我们把投影算符 $\Pi$（投影到规范不变态上）作用于任意截面 $\Psi_{0}[A_{z}]$ 时，可以写成
\begin{equation}
  (\Pi\Psi_{0})[A_{z}]:=\int_{\mathrm{Map}(\Sigma_2\to G)}e^{f^{\pm}(A_{z},g)}\,\Psi_{0}[A^{g}_{z}]
  \tag{3.65}
  \label{eq:3.65}
\end{equation}
我们的投影算符就是二维量子场论——即群 $G$ 的 WZW 理论——的一个路径积分。我们将在第 4 节进一步解释这一理论。

对于 $f^{-}$ 的情形，一个特别有用的检验函数是取
\begin{equation}
  \Psi_{0}[A_{z}]=\exp\left[\mathrm{i}\xi k\frac{1}{2\pi}\int\operatorname{tr}\left(J^{0,1}A^{1,0}\right)\right]
  \tag{3.66}
  \label{eq:3.66}
\end{equation}
其中我们选取某个外部的 $(0,1)$ 形式。于是
\begin{equation}
  (\Pi\Psi_{0})[A_{z}]:=e^{-\mathrm{i}\xi k\frac{1}{2\pi}\int\operatorname{tr}\left(A^{1,0}J^{0,1}\right)}\int_{\mathrm{Map}(\Sigma_2\to G)}e^{\mathrm{i}\xi k\,S^{-}[g;A^{1,0},J^{0,1}]}
  \tag{3.67}
  \label{eq:3.67}
\end{equation}
其中 $S^{-}[g;A,\bar A]$ 是带规范场的 WZW 作用量：
\begin{equation}
  \begin{aligned}
    S^{-}[g;A,\bar A]&:=S^{-}(g)+\\
    &\quad+\frac{1}{2\pi}\int_{\Sigma_2}\left[\operatorname{tr}\left(A\bar\partial g\,g^{-1}+\bar A g^{-1}\partial g-g^{-1}Ag\bar A+A\bar A\right)\right]
  \end{aligned}
  \tag{3.68}
  \label{eq:3.68}
\end{equation}
其中 $A$ 是联络的 $(1,0)$ 部分，$\bar A$ 是联络的 $(0,1)$ 部分。
% [Tutor Note: (3.68) 里 $g^{-1}Ag\bar A$ 与 $A\bar A$ 两项的符号是整条论证中最容易出错的地方。要点是：带规范场的 WZW 作用量必须在 $A\to A^{h}$、$g\to h^{-1}g$（或相应的手征变换）下保持不变，而 (3.66) 选取的 $\Psi_{0}$ 恰好把外部 $(0,1)$ 形式 $J^{0,1}$ 与 $A^{1,0}$ 配对；用这一不变性可以逐项定出 $+A\bar A$、$-g^{-1}Ag\bar A$ 的相对符号。若把它写成 $-A\bar A$，则 (3.67) 的规范不变性会破坏。]

{\small\itshape ［原文标记：疑难：用 $f^{+}$ 无法干净地组合成带规范场的 WZW 作用量。为什么不行？］}
% [原文待办标记，原文 p.95，位于页边]

\textbf{评论}

\begin{enumerate}
  \item 若 $G$ 具有酉结构，即 $g^{\dagger}=g^{-1}$（例如 $\mathfrak{g}\subset\mathfrak{gl}(n,\mathbb{C})$ 是反 Hermite 矩阵的 Lie 子代数），则按我们的约定
\begin{equation}
  \mathrm{i}\int_{\Sigma_2}\operatorname{tr}\left(g^{-1}\partial g\,g^{-1}\bar\partial g\right)
  \tag{3.69}
  \label{eq:3.69}
\end{equation}
  是负定的。因此，若我们想要 \eqref{eq:3.65} 中关于映射 $g:\Sigma_2\to G$ 的路径积分在形式上收敛，就应当在 $\xi k<0$ 时选 $f^{+}$，在 $\xi k>0$ 时选 $f^{-}$。
\end{enumerate}
